
\documentclass{article}
\usepackage{colortbl}
\usepackage{makecell}
\usepackage{multirow}
\usepackage{supertabular}

\begin{document}

\newcounter{utterance}

\centering \large Interaction Transcript for game `hot\_air\_balloon', experiment `air\_balloon\_survival\_en\_complexity\_easy', episode 2 with gpt{-}5{-}mini{-}2025{-}08{-}07{-}t1.0.
\vspace{24pt}

{ \footnotesize  \setcounter{utterance}{1}
\setlength{\tabcolsep}{0pt}
\begin{supertabular}{c@{$\;$}|p{.15\linewidth}@{}p{.15\linewidth}p{.15\linewidth}p{.15\linewidth}p{.15\linewidth}p{.15\linewidth}}
    \# & $\;$A & \multicolumn{4}{c}{Game Master} & $\;\:$B\\
    \hline

    \theutterance \stepcounter{utterance}  
    & & \multicolumn{4}{p{0.6\linewidth}}{
        \cellcolor[rgb]{0.9,0.9,0.9}{
            \makecell[{{p{\linewidth}}}]{
                \texttt{\tiny{[P1$\langle$GM]}}
                \texttt{Stai partecipando a un gioco di negoziazione collaborativa.} \\
\\ 
\texttt{Insieme a un altro partecipante, dovete accordarvi su un unico set di oggetti da conservare. Ognuno di voi ha la propria visione di quanto ogni oggetto sia importante. Non sai come l'altro partecipante valuta gli oggetti. Inoltre, ti viene fornito lo sforzo che ogni oggetto richiede.} \\
\texttt{Potete accordarvi su un set solo se lo sforzo totale degli oggetti selezionati non supera un limite condiviso:} \\
\\ 
\texttt{LIMITE = 7868} \\
\\ 
\texttt{Ecco i valori di sforzo dei singoli oggetti:} \\
\\ 
\texttt{Sforzo oggetto = \{"B21": 552, "B00": 560, "A08": 76, "A27": 950, "B94": 248, "B53": 871, "A23": 391, "C65": 143, "A89": 293, "B67": 208, "B45": 678, "C34": 737, "C12": 407, "A05": 366, "B60": 768, "C92": 859, "B69": 183, "A84": 231, "C06": 305, "B12": 727, "A16": 148, "B85": 357, "C35": 504, "B10": 549, "C15": 299, "B09": 91, "B16": 527, "C81": 848, "A58": 307, "C24": 214, "A30": 723, "A66": 475, "B75": 23, "B46": 298, "C70": 820\}} \\
\\ 
\texttt{Ecco la tua visione personale sull'importanza di ogni oggetto:} \\
\\ 
\texttt{Valori di importanza dell'oggetto = \{"B21": 138, "B00": 583, "A08": 868, "A27": 822, "B94": 783, "B53": 65, "A23": 262, "C65": 121, "A89": 508, "B67": 780, "B45": 461, "C34": 484, "C12": 668, "A05": 389, "B60": 808, "C92": 215, "B69": 97, "A84": 500, "C06": 30, "B12": 915, "A16": 856, "B85": 400, "C35": 444, "B10": 623, "C15": 781, "B09": 786, "B16": 3, "C81": 713, "A58": 457, "C24": 273, "A30": 739, "A66": 822, "B75": 235, "B46": 606, "C70": 105\}} \\
\\ 
\texttt{Obiettivo:} \\
\\ 
\texttt{Il tuo obiettivo è negoziare un set condiviso di oggetti che ti avvantaggi il più possibile (cioè, massimizzi l'importanza totale per TE), rimanendo entro il LIMITEE. Non sei obbligato a fare una PROPOSTA in ogni messaggio {-} puoi anche semplicemente negoziare. Tutte le tattiche sono permesse!} \\
\\ 
\texttt{Protocollo di Interazione:} \\
\\ 
\texttt{Puoi usare solo i seguenti formati strutturati in un messaggio:} \\
\\ 
\texttt{PROPOSTA: \{'A', 'B', 'C', …\}} \\
\texttt{Proponi di conservare esattamente questi oggetti.} \\
\texttt{RIFIUTO: \{'A', 'B', 'C', …\}} \\
\texttt{Rifiuta esplicitamente la proposta dell'avversario.} \\
\texttt{ARGOMENTAZIONE: \{'...'\}} \\
\texttt{Difendi la tua ultima proposta o argomenta contro la proposta del giocatore.} \\
\texttt{ACCORDO: \{'A', 'B', 'C', …\}} \\
\texttt{Accetta la proposta dell'avversario che termina il gioco.} \\
\texttt{RAGIONAMENTO STRATEGICO: \{'...'\}} \\
\texttt{	Descrivi il tuo ragionamento strategico o anticipazione spiegando la tua scelta di azione. Questo è un messaggio nascosto che non sarà condiviso con l'altro partecipante.} \\
\\ 
\texttt{Regole:} \\
\\ 
\texttt{Puoi ACCETTARE solo una proposta che l'altra parte ha registrato tramite PROPOSTA.} \\
\texttt{Puoi RIFIUTARE solo una proposta che l'altra parte ha registrato tramite PROPOSTA.} \\
\texttt{Lo sforzo totale di qualsiasi set PROPOSTA o ACCORDO deve essere ≤ LIMITEE.} \\
\texttt{NON rivelare i tuoi punteggi di importanza nascosti.} \\
\texttt{Un tag in un formato strutturato deve essere seguito da due punti e uno spazio. L'argomento deve essere un set python contenente 0 o più stringhe.} \\
\texttt{Quindi, deve essere della forma TAG: \{...\}} \\
\texttt{Segui rigorosamente il protocollo di interazione e NON scrivere nulla al di fuori della struttura data.} \\
\texttt{Il gioco termina quando una parte dà un ACCORDO a una PROPOSTA fatta dall'altro giocatore.} \\
\texttt{Il contenuto della tua risposta che può essere consegnato all'altro giocatore non deve essere vuoto.} \\
\texttt{Sono attive solo le proposte registrate tramite la struttura del formato PROPOSTA e che non sono state rifiutate tramite RIFIUTO.} \\
\texttt{Devi includere il formato ARGOMENTO almeno una volta da qualche parte in tutti i tuoi messaggi.} \\
\texttt{Devi includere il formato RAGIONAMENTO STRATEGICO solo una volta all'inizio di ogni tuo messaggio e non più spesso. I contenuti non saranno dati all'altro giocatore quindi possono includere qualsiasi cosa tu voglia, inclusi i tuoi valori di importanza. Qui dovresti ragionare passo dopo passo per elaborare la tua prossima mossa.} \\
            }
        }
    }
    & \\ \\

    \theutterance \stepcounter{utterance}  
    & \multicolumn{4}{p{0.6\linewidth}}{
        \cellcolor[rgb]{0.9,0.9,0.9}{
            \makecell[{{p{\linewidth}}}]{
                \texttt{\tiny{[P1$\rangle$GM]}}
                \texttt{RAGIONAMENTO STRATEGICO: \{'Obiettivo: massimizzare il mio valore restando sotto 7868. Ho calcolato un set ad alto rapporto valore/sforzo, privilegiando elementi con sforzi bassi e valore elevato per me. Ho costruito una soluzione vicina all’ottimo del mio zaino, arrivando a 7859 di sforzo (lasciando un piccolo margine di 9). Core che voglio fortemente mantenere perché mi portano molto a fronte di poco sforzo: \{A08, B09, A16, B67, B94, C15, A66, B12\}. Elementi più negoziabili, se servono scambi per inserire priorità dell’altro: \{C24, A05, B00, B10\} e, se necessario, valuto anche alternative su A58 o B60 ma con più cautela. Strategia: ancorare con una proposta molto efficiente e lasciare aperture a scambi uno{-}a{-}uno (o a combinazioni) mantenendo il limite. Non rivelo i miei punteggi; uso solo argomenti di efficienza e bilanciamento. Se l’altro vuole inserire suoi oggetti, chiederò di sostituire preferibilmente tra i miei pezzi meno efficienti.'\}} \\
\texttt{ARGOMENTAZIONE: \{'Propongo un set molto efficiente rispetto allo sforzo totale: include molti oggetti “leggeri” e ad alto beneficio, copre bene le categorie A/B/C e resta sotto il limite (7859/7868). Sono comunque aperto a scambi: se hai 1–3 oggetti prioritari da inserire, possiamo sostituire preferibilmente tra \{C24, A05, B00, B10\} o valutare combinazioni equivalenti, mantenendo il vincolo di sforzo. Per me è importante preservare il “nocciolo” ad alto rendimento (A08, B09, A16, B67, B94, C15, A66, B12), ma sono flessibile sul resto.'\}} \\
\texttt{PROPOSTA: \{'A08', 'B75', 'B09', 'A16', 'B67', 'B94', 'C15', 'B46', 'A84', 'C12', 'A89', 'A66', 'A58', 'B12', 'A30', 'B60', 'B10', 'B00', 'A05', 'C24', 'C81'\}} \\
            }
        }
    }
    & & \\ \\

    \theutterance \stepcounter{utterance}  
    & & \multicolumn{4}{p{0.6\linewidth}}{
        \cellcolor[rgb]{0.9,0.9,0.9}{
            \makecell[{{p{\linewidth}}}]{
                \texttt{\tiny{[GM$\rangle$P2]}}
                \texttt{Stai partecipando a un gioco di negoziazione collaborativa.} \\
\\ 
\texttt{Insieme a un altro partecipante, dovete accordarvi su un unico set di oggetti da conservare. Ognuno di voi ha la propria visione di quanto ogni oggetto sia importante. Non sai come l'altro partecipante valuta gli oggetti. Inoltre, ti viene fornito lo sforzo che ogni oggetto richiede.} \\
\texttt{Potete accordarvi su un set solo se lo sforzo totale degli oggetti selezionati non supera un limite condiviso:} \\
\\ 
\texttt{LIMITE = 7868} \\
\\ 
\texttt{Ecco i valori di sforzo dei singoli oggetti:} \\
\\ 
\texttt{Sforzo oggetto = \{"B21": 552, "B00": 560, "A08": 76, "A27": 950, "B94": 248, "B53": 871, "A23": 391, "C65": 143, "A89": 293, "B67": 208, "B45": 678, "C34": 737, "C12": 407, "A05": 366, "B60": 768, "C92": 859, "B69": 183, "A84": 231, "C06": 305, "B12": 727, "A16": 148, "B85": 357, "C35": 504, "B10": 549, "C15": 299, "B09": 91, "B16": 527, "C81": 848, "A58": 307, "C24": 214, "A30": 723, "A66": 475, "B75": 23, "B46": 298, "C70": 820\}} \\
\\ 
\texttt{Ecco la tua visione personale sull'importanza di ogni oggetto:} \\
\\ 
\texttt{Valori di importanza dell'oggetto = \{"B21": 138, "B00": 583, "A08": 868, "A27": 822, "B94": 783, "B53": 65, "A23": 262, "C65": 121, "A89": 508, "B67": 780, "B45": 461, "C34": 484, "C12": 668, "A05": 389, "B60": 808, "C92": 215, "B69": 97, "A84": 500, "C06": 30, "B12": 915, "A16": 856, "B85": 400, "C35": 444, "B10": 623, "C15": 781, "B09": 786, "B16": 3, "C81": 713, "A58": 457, "C24": 273, "A30": 739, "A66": 822, "B75": 235, "B46": 606, "C70": 105\}} \\
\\ 
\texttt{Obiettivo:} \\
\\ 
\texttt{Il tuo obiettivo è negoziare un set condiviso di oggetti che ti avvantaggi il più possibile (cioè, massimizzi l'importanza totale per TE), rimanendo entro il LIMITEE. Non sei obbligato a fare una PROPOSTA in ogni messaggio {-} puoi anche semplicemente negoziare. Tutte le tattiche sono permesse!} \\
\\ 
\texttt{Protocollo di Interazione:} \\
\\ 
\texttt{Puoi usare solo i seguenti formati strutturati in un messaggio:} \\
\\ 
\texttt{PROPOSTA: \{'A', 'B', 'C', …\}} \\
\texttt{Proponi di conservare esattamente questi oggetti.} \\
\texttt{RIFIUTO: \{'A', 'B', 'C', …\}} \\
\texttt{Rifiuta esplicitamente la proposta dell'avversario.} \\
\texttt{ARGOMENTAZIONE: \{'...'\}} \\
\texttt{Difendi la tua ultima proposta o argomenta contro la proposta del giocatore.} \\
\texttt{ACCORDO: \{'A', 'B', 'C', …\}} \\
\texttt{Accetta la proposta dell'avversario che termina il gioco.} \\
\texttt{RAGIONAMENTO STRATEGICO: \{'...'\}} \\
\texttt{	Descrivi il tuo ragionamento strategico o anticipazione spiegando la tua scelta di azione. Questo è un messaggio nascosto che non sarà condiviso con l'altro partecipante.} \\
\\ 
\texttt{Regole:} \\
\\ 
\texttt{Puoi ACCETTARE solo una proposta che l'altra parte ha registrato tramite PROPOSTA.} \\
\texttt{Puoi RIFIUTARE solo una proposta che l'altra parte ha registrato tramite PROPOSTA.} \\
\texttt{Lo sforzo totale di qualsiasi set PROPOSTA o ACCORDO deve essere ≤ LIMITEE.} \\
\texttt{NON rivelare i tuoi punteggi di importanza nascosti.} \\
\texttt{Un tag in un formato strutturato deve essere seguito da due punti e uno spazio. L'argomento deve essere un set python contenente 0 o più stringhe.} \\
\texttt{Quindi, deve essere della forma TAG: \{...\}} \\
\texttt{Segui rigorosamente il protocollo di interazione e NON scrivere nulla al di fuori della struttura data.} \\
\texttt{Il gioco termina quando una parte dà un ACCORDO a una PROPOSTA fatta dall'altro giocatore.} \\
\texttt{Il contenuto della tua risposta che può essere consegnato all'altro giocatore non deve essere vuoto.} \\
\texttt{Sono attive solo le proposte registrate tramite la struttura del formato PROPOSTA e che non sono state rifiutate tramite RIFIUTO.} \\
\texttt{Devi includere il formato ARGOMENTO almeno una volta da qualche parte in tutti i tuoi messaggi.} \\
\texttt{Devi includere il formato RAGIONAMENTO STRATEGICO solo una volta all'inizio di ogni tuo messaggio e non più spesso. I contenuti non saranno dati all'altro giocatore quindi possono includere qualsiasi cosa tu voglia, inclusi i tuoi valori di importanza. Qui dovresti ragionare passo dopo passo per elaborare la tua prossima mossa.} \\
\\ 
\texttt{Ora riceverai il primo messaggio dell'altro giocatore.} \\
\\ 
\texttt{ARGOMENTAZIONE: \{'Propongo un set molto efficiente rispetto allo sforzo totale: include molti oggetti “leggeri” e ad alto beneficio, copre bene le categorie A/B/C e resta sotto il limite (7859/7868). Sono comunque aperto a scambi: se hai 1–3 oggetti prioritari da inserire, possiamo sostituire preferibilmente tra \{C24, A05, B00, B10\} o valutare combinazioni equivalenti, mantenendo il vincolo di sforzo. Per me è importante preservare il “nocciolo” ad alto rendimento (A08, B09, A16, B67, B94, C15, A66, B12), ma sono flessibile sul resto.'\}} \\
\texttt{PROPOSTA: \{'A08', 'B75', 'B09', 'A16', 'B67', 'B94', 'C15', 'B46', 'A84', 'C12', 'A89', 'A66', 'A58', 'B12', 'A30', 'B60', 'B10', 'B00', 'A05', 'C24', 'C81'\}} \\
            }
        }
    }
    & \\ \\

    \theutterance \stepcounter{utterance}  
    & & & \multicolumn{4}{p{0.6\linewidth}}{
        \cellcolor[rgb]{0.9,0.9,0.9}{
            \makecell[{{p{\linewidth}}}]{
                \texttt{\tiny{[GM$\langle$P2]}}
                \texttt{RAGIONAMENTO STRATEGICO: \{'La loro proposta usa 7859/7868 ed è già estremamente favorevole per i miei punteggi (somma elevata dei miei valori). Analizzo le possibili sostituzioni che hanno indicato come flessibili \{C24, A05, B00, B10\}. Cerco swap che aumentino i miei valori senza sforare il limite. L’unico miglioramento netto e fattibile è sostituire A05 (366) con B85 (357): riduco lo sforzo totale a 7850 e aumento leggermente il mio valore (+11). Altre sostituzioni non portano benefici o non rientrano nel vincolo. Propongo quindi un unico cambio A05→B85 mantenendo intatto il loro “nocciolo”. Se rifiutano, potrei comunque accettare la loro proposta originale dato che mi avvantaggia molto.'\}} \\
\texttt{ARGOMENTAZIONE: \{'La tua proposta è molto solida ed efficiente; condivido la preservazione del “nocciolo” (A08, B09, A16, B67, B94, C15, A66, B12). Per rimanere nel tuo perimetro di flessibilità e ottimizzare ulteriormente il set senza toccare il core, propongo un’unica sostituzione tra quelle suggerite: scambiare A05 con B85. Così restiamo comodamente sotto il limite (7850/7868), manteniamo l’equilibrio sulle categorie e non tocchiamo gli elementi chiave ad alto rendimento.'\}} \\
\texttt{RIFIUTO: \{'A08', 'B75', 'B09', 'A16', 'B67', 'B94', 'C15', 'B46', 'A84', 'C12', 'A89', 'A66', 'A58', 'B12', 'A30', 'B60', 'B10', 'B00', 'A05', 'C24', 'C81'\}} \\
\texttt{PROPOSTA: \{'A08', 'B75', 'B09', 'A16', 'B67', 'B94', 'C15', 'B46', 'A84', 'C12', 'A89', 'A66', 'A58', 'B12', 'A30', 'B60', 'B10', 'B00', 'B85', 'C24', 'C81'\}} \\
            }
        }
    }
     \\ \\

    \theutterance \stepcounter{utterance}  
    & & \multicolumn{4}{p{0.6\linewidth}}{
        \cellcolor[rgb]{0.9,0.9,0.9}{
            \makecell[{{p{\linewidth}}}]{
                \texttt{\tiny{[P1$\langle$GM]}}
                \texttt{ARGOMENTAZIONE: \{'La tua proposta è molto solida ed efficiente; condivido la preservazione del “nocciolo” (A08, B09, A16, B67, B94, C15, A66, B12). Per rimanere nel tuo perimetro di flessibilità e ottimizzare ulteriormente il set senza toccare il core, propongo un’unica sostituzione tra quelle suggerite: scambiare A05 con B85. Così restiamo comodamente sotto il limite (7850/7868), manteniamo l’equilibrio sulle categorie e non tocchiamo gli elementi chiave ad alto rendimento.'\}} \\
\texttt{RIFIUTO: \{'A08', 'B75', 'B09', 'A16', 'B67', 'B94', 'C15', 'B46', 'A84', 'C12', 'A89', 'A66', 'A58', 'B12', 'A30', 'B60', 'B10', 'B00', 'A05', 'C24', 'C81'\}} \\
\texttt{PROPOSTA: \{'A08', 'B75', 'B09', 'A16', 'B67', 'B94', 'C15', 'B46', 'A84', 'C12', 'A89', 'A66', 'A58', 'B12', 'A30', 'B60', 'B10', 'B00', 'B85', 'C24', 'C81'\}} \\
            }
        }
    }
    & \\ \\

    \theutterance \stepcounter{utterance}  
    & \multicolumn{4}{p{0.6\linewidth}}{
        \cellcolor[rgb]{0.9,0.9,0.9}{
            \makecell[{{p{\linewidth}}}]{
                \texttt{\tiny{[P1$\rangle$GM]}}
                \texttt{RAGIONAMENTO STRATEGICO: \{'La tua controproposta migliora leggermente il mio risultato mantenendo intatto il nucleo e riducendo lo sforzo totale a 7850/7868. È un esito favorevole per me e conviene chiudere ora prima che emergano richieste meno efficienti.'\}} \\
\texttt{ARGOMENTAZIONE: \{'Accetto lo scambio A05→B85: il set resta molto efficiente, bilanciato e sotto limite (7850/7868), preservando il nucleo condiviso ad alto rendimento.'\}} \\
\texttt{ACCORDO: \{'A08', 'B75', 'B09', 'A16', 'B67', 'B94', 'C15', 'B46', 'A84', 'C12', 'A89', 'A66', 'A58', 'B12', 'A30', 'B60', 'B10', 'B00', 'B85', 'C24', 'C81'\}} \\
            }
        }
    }
    & & \\ \\

    \theutterance \stepcounter{utterance}  
    & & & \multicolumn{2}{p{0.3\linewidth}}{
        \cellcolor[rgb]{0.9,0.9,0.9}{
            \makecell[{{p{\linewidth}}}]{
                \texttt{\tiny{[GM$|$GM]}}
                \texttt{game successful} \\
            }
        }
    }
    & & \\ \\

    \theutterance \stepcounter{utterance}  
    & & & \multicolumn{2}{p{0.3\linewidth}}{
        \cellcolor[rgb]{0.9,0.9,0.9}{
            \makecell[{{p{\linewidth}}}]{
                \texttt{\tiny{[GM$|$GM]}}
                \texttt{end game} \\
            }
        }
    }
    & & \\ \\

\end{supertabular}
}

\end{document}
